\chapter*{Kết luận và hướng phát triển}
\addcontentsline{toc}{chapter}{Kết luận và hướng phát triển}
Đồ án đã xây dựng một quy trình từ dữ liệu Patent/Publication đến khám phá chủ đề, truy hồi và kiểm chứng trên 13.702 tài liệu có truy vết. LDA $K=20$ cung cấp phân bố chủ đề xác suất; BGE và BERTopic bổ sung biểu diễn ngữ nghĩa và các cụm chi tiết. Fine-tuning block cuối được thực hiện thực với checkpoint chọn bằng validation; cải thiện proxy Title--Abstract được ghi nhận trong điều kiện benchmark đã chốt.

Nhánh full-Abstract chunking tạo 29.139 đoạn và bảo đảm bao phủ token, đồng thời làm thay đổi cấu trúc topic: coherence/diversity tăng, noise cũng tăng. Kết quả không được diễn giải thành chiến thắng tổng thể. Search v2 triển khai lexical/dense retrieval, RRF và cross-encoder, có kiểm soát input, lọc metadata và trả passage. Pretrained là mặc định thiết kế; fine-tuned/chunked và SE1/SE2/SE3 được giữ để đối chứng thay vì loại bỏ các kết quả bất lợi.

Đóng góp đánh giá là tách relevance, topic metrics và source recovery; xây dựng pool/qrels có provenance, input validation và uncertainty theo đúng đơn vị lấy mẫu. Kết quả relevance AI còn một phần, chưa có nhãn chuyên gia độc lập; synthetic mới đủ 197/200 query và tập source chủ yếu thuộc train. Do đó, báo cáo chưa xác lập phương pháp tốt nhất hoặc mức độ chính xác tìm kiếm chung. Đây là hệ thống tích hợp đã có kiểm chứng kỹ thuật và bằng chứng thực nghiệm có điều kiện, không phải tuyên bố thuật toán mới.

Ưu tiên tiếp theo là tiếp tục đúng checkpoint khi quota API cho phép, hoàn tất A/B/C còn thiếu và giữ nguyên query/rankings đã đóng băng. Sau đó cần lấy mẫu đối chiếu chuyên gia với rubric rõ, phân tích mức agreement và lỗi có bằng chứng trước khi lựa chọn mô hình. Việc thêm nhãn chuyên gia phải ghi provenance thật, không đổi nhãn AI thành nhãn người.

Hướng nghiên cứu tiếp theo là mở rộng truy vấn từ người dùng độc lập, tách development/test theo nhóm ý định và đo lỗi theo hệ pin, công thức, điều kiện, quan hệ. Nếu sửa preprocessing hoặc thử fine-tuning mới, cần một version riêng, dữ liệu query--positive--hard-negative đã rà soát, test giữ kín theo nhóm và nhiều seed. Ablation phải thay từng yếu tố Title context, max length, full coverage và pooling để quy nguyên nhân. Cải thiện vận hành cần workload lớn hơn, cold start, RAM và concurrency; không tối ưu chỉ để làm điểm trên hai query hiện tại cao hơn.
